\section{Regularity of the strip partition function}\label{c:smoothness-section}

The renewal identity alone gives a tail estimate for the height of one
irreducible bridge.  Prescribing the terminal port requires an estimate
for a single height.  Our first step is to show that the strip partition
function varies smoothly on its natural scale: every fixed forward
difference gains the corresponding power of the strip width.

\begin{proposition}\label{c:strip-smoothness}
Let $B_N$ be the critical mass of strict bridges from one fixed port to
all ports on the opposite pure cut at integral height $N$.  With
$\Delta f_N=f_{N+1}-f_N$, for every fixed integer $\ell\ge0$ and every
$\eta>0$ there is $C_{\ell,\eta}<\infty$ such that
\[
 |\Delta^\ell B_N|\le C_{\ell,\eta}N^{-1/4-\ell+\eta}
 \qquad(N\ge1).
\]
\end{proposition}

The probabilistic application needs precisely Proposition~\ref{c:strip-smoothness}:
its passage to an exact height atom begins in Section~\ref{c:atoms-section}.
We now prove the proposition from the exact finite strip representation and its positive
interpolation spaces.  We state this input before introducing the
holomorphic kernels used to estimate it.  The spectral parameter
$\lambda=\pi/8$ below is unrelated to a path length.
There are three analytic tasks.  We first transport the residual space
to the disk and control its evaluation kernel, including the boundary
column for the non-square-integrable input.  We then recover the residuals at $0$ and $-4$
from their Fourier transforms, retaining every fixed difference in $N$.
Finally a separate positivity argument bounds the denominator away from
zero; only then do we estimate the quotient for $B_N$.

\subsection{The finite strip input}

Put $b=3\lambda$, $C_t(x)=\cosh(t\sqrt x)$,
$\mathcal D(x)=2C_\lambda(x)-1$, and
$a_*=\mathcal D(-4)=2\cos(2\lambda)-1>0$.
Define the positive measure $\mu$ on $(0,\infty)$ by
\begin{equation}\label{c:strip-measure}
 d\mu(x)=\frac{z\mathcal D(z^2)}{\sinh(4\lambda z)}\,dz,
 \qquad x=z^2,\quad z>0.
\end{equation}
For $N\ge1$, set
\[
 \begin{split}
 M_N&=\operatorname{span}\{1,x,\ldots,x^{N-1},
              \partial_t^j C_t|_{t=\lambda}:0\le j<N\},\\
 c_1&=C_{\lambda/2},\quad c_3=C_{3\lambda/2},\quad
 h_0=\mathcal D^{-1},\quad h_1=c_1/\mathcal D,\\
 g&=\int_0^\lambda C_t\,dt,
 \quad F_{0,N}=M_N+\mathbb R c_1,
 \quad G_N=M_N+\mathbb R g.
 \end{split}
\]
Let $\Pi_Nh\in F_{0,N}$ be the unique oblique projection preserving
all pairings against $G_N$ in $L^2(\mu)$, and write
$r_Nh=h-\Pi_Nh$.  The projection is also defined for $c_3$, whose
pairings with $G_N$ are finite although $c_3\notin L^2(\mu)$.

We use the following finite facts from the strip analysis in
\cite[Theorem 2.1 and Lemmas 3.1--3.4]{UniformMarked}.
\begin{enumerate}
\item The above projection exists.  Whenever the denominator is nonzero,
\begin{equation}\label{c:strip-residual-formula}
 B_N=2a_*\left[
 r_Nh_0(-4)-
 \frac{r_N(c_3-h_1)(-4)\,r_Nh_0(0)}{r_N(c_3-h_1)(0)}
 \right].
\end{equation}
The corresponding undivided statement is that there is a scalar $s_N$ such that
$2B_N=4a_*[r_Nh_0(-4)+s_Nr_Nh_0(0)]$ and
$r_N(c_3-h_1)(-4)+s_Nr_N(c_3-h_1)(0)=0$.

\item For $0<t<\lambda$, put $G_{N,t}=M_N+\mathbb R C_t$.
On ordered positive nodes $X=(x_1<\cdots<x_{2N+1})$, the product of
evaluation determinants for $(F_{0,N},G_{N,t})$, times a positive
base weight at every node, defines a positive density whenever its
moments are finite.  Its mean interpolant is the oblique projection
against $G_{N,t}$.  If $e_0c_3$ denotes the interpolation residual at
zero, then $-e_0c_3$ is positive and increasing in each ordered node.
Multiplying the base weight by an increasing positive function increases
its expectation.  Replacing $G_{N,t}$ by $G_N$ is a positive mixture
of these densities.

\item Let $F_N=F_{0,N}+\mathbb R c_3$,
$A_N=\{f\in F_N:f(0)=0\}$, and
$U_N=\{f\in M_N:f(0)=0\}$.  For either determinant-product density
$(F_{0,N},G_{N,t})$ or $(A_N,G_{N,t})$, expectations of
$\prod_i r(x_i)$, with $r(x)=(x+\alpha)/(x+\beta)$ and fixed
$\alpha,\beta>0$, differ by bounded factors from the corresponding
symmetric expectations on $(M_N,M_N)$ or $(U_N,U_N)$.
The factors are uniform in $N,t$ and remain so after any fixed number
of such rational changes of the base weight. This is the interlacing
reduction in the proof of \cite[Lemma 3.2]{UniformMarked}: its positive
weight and exponential-domination hypotheses persist under each such
change, and the bound uses only the range of $\log r$, not $t$.  The two symmetric
expectations themselves differ by a bounded factor.  These statements
follow from interlacing of the positive roots of the compressed
multiplication forms; a codimension-one compression changes the
logarithmic product by at most the range of $\log r$.

\item For the probability measure $\mu_c$ proportional to
$dz/\cosh(cz)$, $x=z^2$, let
$\alpha_c=\pi/(2c)$ and $w_t=\tan^2(\alpha_ct)$.
Then
\begin{equation}\label{c:polynomial-generator}
 A_t(x)=\cos(\alpha_ct)C_t(x),\qquad
 \int A_tA_s\,d\mu_c=\frac1{1-w_tw_s}.
\end{equation}
Thus $A_t=\sum_{j\ge0}P_j(x)w_t^j$ generates orthonormal polynomials
$P_j$ of degree $j$.  The identity and its differentiated forms hold
when $|t|+|s|<c$.
\end{enumerate}
All base weights below have the required exponential domination.  In
particular these inputs apply to $\mu$, to
$dW=z\,dz/\sinh(bz)$, and to $r^jW$ for fixed $j$ and fixed positive
rational function $r$ of the indicated form.  We will prove both the estimates and the nonvanishing
needed to use \eqref{c:strip-residual-formula}.

\subsection{Transforming the residual space}

Set
\[
 v(\omega)=\tan\frac{\pi\omega}{2b},\qquad
 J(\omega)=v'(\omega),\qquad
 Q(v)=\frac{v^2(v^2-v(\lambda)^2)}{1-v^2v(\lambda)^2}.
\]
The strip $|\operatorname{Re}\omega|<b/2$ maps bijectively to the unit
disk.  The rational function $Q$ is an even finite Blaschke product;
its zeros are $0$ with multiplicity two and $\pm v(\lambda)$ with
multiplicity one.  For an even holomorphic function $A$ on the disk,
define
\[
 (T_NA)(\omega)=J(\omega)Q(v(\omega))^N A(v(\omega)).
\]
The factor $Q^N$ encodes exactly the orthogonality conditions against
$M_N$.

Indeed, the cosine transform identifies $h\in M_N^\perp\subset
L^2(\mu)$ with
\begin{equation}\label{c:residual-transform}
 \int h(x)C_\omega(x)\,d\mu(x)=T_NA(\omega).
\end{equation}
The transform is even, has a zero of order $2N$ at zero, and zeros of
order $N$ at $\pm\lambda$.  Dividing by $JQ^N$ therefore gives an
even holomorphic $A$ on the disk.  Conversely these zero conditions
imply the required orthogonality.

Here is the Hilbert-space content of this identification.  Let
$w_\mu(z)$ and $w_W(z)$ be the even extensions of the $z$-densities of
$\mu$ and $W$.  The Fourier datum of the transform $F=T_NA$ is
$f(z)=h(z^2)w_\mu(z)$, so
\[
 F(\omega)=\frac12\int_{\mathbb R}f(z)e^{\omega z}\,dz,
 \qquad
 \|h\|_{L^2(\mu)}^2
 =\frac12\int_{\mathbb R}\frac{|f(z)|^2}{w_\mu(z)}\,dz.
\]
We compare norms while holding $F$, and hence $f$, fixed.  In the
comparison measure $W$, the same transform is represented by
$h_W=h\,d\mu/dW$; its squared norm replaces $w_\mu$ in the last
integral by $w_W$.  Moving to a parallel line multiplies the Fourier
datum by $e^{z\operatorname{Re}\omega}$.  Integrating Plancherel's
identity across the strip supplies the reciprocal weight
$\sinh(bz)/z=1/w_W(z)$.  The conformal Jacobian $J$ then gives the
transported comparison norm
\begin{equation}\label{c:model-norm}
 \mathcal G_N^0(A,A)
 =c_b\int_{|v|<1}|Q(v)|^{2N}|A(v)|^2\,d^2v,
 \qquad c_b>0.
\end{equation}
For completeness, the converse Fourier representation follows by
Fourier transforming the Cauchy--Riemann equation locally in frequency:
the data on parallel lines must have the same exponential multiplier.
Strip square integrability gives the stated weighted $L^2$ data, and
division by the even weight gives $h$.  Thus no unproved surjectivity
of the transform is being used.

Write $\mathcal G_N(A,A)=\|h\|_{L^2(\mu)}^2$ for the norm in
\eqref{c:residual-transform}.  The density ratio of $\mu$ to $W$ is
bounded above and below by positive constants, so
$\mathcal G_N\asymp\mathcal G_N^0$ uniformly in $N$.
More precisely, their difference is a fixed bounded sesquilinear form
on the two traces $T_NA(\pm d_{\mathrm{tr}}+iy)$, for a fixed $d_{\mathrm{tr}}<b/2$ sufficiently
close to $b/2$.  To see this, the difference of the reciprocal Fourier
weights is, up to a fixed factor,
\[
 z^{-1}\left\{\frac{\sinh(4\lambda z)}{\mathcal D(z^2)}
                    -\sinh(bz)\right\}.
\]
It is bounded at zero and is $O(e^{2d_{\mathrm{tr}}|z|})$ at infinity: the terms of
largest exponential order cancel.  Plancherel on the two parallel
lines therefore bounds the difference form by their $L^2$ trace norms.

If $A(v)=\sum_{j\ge0}a_jv^{2j}$, both transported norms are equivalent
to
\begin{equation}\label{c:coefficient-norm}
 \|A\|_N^2=\sum_{j\ge0}\frac{|a_j|^2}{N+j}.
\end{equation}
The boundary phase speed of the finite Blaschke product $Q$ is smooth
and strictly positive.  Hence $|Q(re^{i\theta})|^{2N}$ is bounded
above by $e^{-cN(1-r)}$, and bounded below by
$e^{-CN(1-r)}$ for $r$ near one.  Angular orthogonality of the powers,
followed by radial integration, proves the equivalence.

Every fixed $\ell$th forward difference of $\mathcal G_N^0$ costs
$O_\ell(N^{-\ell})$ as a form in \eqref{c:coefficient-norm}, since
it inserts $(|Q|^2-1)^\ell$ in the area integral.  The same is true
for $\mathcal G_N$.  On each of the fixed interior traces,
$v\to\pm i$ nontangentially at the two ends and $Q(\pm i)=1$; thus
$|1-Q|\le C(1-|Q|)$.  Each difference of $T_N$ inserts $Q-1$,
and comparison with $T_{\lfloor N/2\rfloor}$ costs $O(N^{-1})$.
Interior trace norms are bounded by strip area norms, by the mean-value
inequality for holomorphic functions.  The discrete product rule
therefore gives the same bound for the correction form.  All norms
with finitely shifted indices are equivalent.  Differencing
$\mathcal G_N^{-1}\mathcal G_N=I$ consequently shows that every
$\ell$th difference of the inverse form costs $O_\ell(N^{-\ell})$
as well.

\subsection{The model evaluation kernel}

Denote by $K_N^0(v,u)$ the evaluation kernel for
$\mathcal G_N^0$, analytic in $v$ and anti-analytic in $u$.
Thus, with the inner product linear in its first argument,
$A(u)=\mathcal G_N^0(A,K_N^0(\cdot,u))$ for every even function in the
space and every $|u|<1$. The following bounds are needed also at boundary
points, provided the two variables do not meet on the even diagonal.

\begin{lemma}\label{c:model-kernel}
On every compact subset of the closed bidisk on which
$|1-v^2\overline u^2|>0$, the kernel has continuous boundary values and
\[
 |\Delta^\ell K_N^0(v,u)|\le C_\ell N^{1-\ell},\qquad
 \frac{K_N^0(v,u)}N\longrightarrow
 \frac{O(v)\overline{O(u)}}{1-v^2\overline u^2}.
\]
Here $O$ is even, holomorphic and nonvanishing on a neighborhood of the
closed disk, and positive on both the real and imaginary diameters.
\end{lemma}

\begin{proof}
Let $P_{jk}$ be the Gram matrix of $v^{2j}$ in
\eqref{c:model-norm}, and $R=P^{-1}$.  Thus
$K_N^0(v,u)=\sum_{j,k\ge0}v^{2j}R_{jk}\overline u^{2k}$ in the
open bidisk.  Write $\delta$ for simultaneous forward difference in
$j,k$.  For all fixed nonnegative integers $\ell,s$,
\begin{equation}\label{c:gram-differences}
 |\Delta^\ell\delta^sP_{jk}|
 \le C_{\ell,s}e^{-c|j-k|}(N+j+k)^{-1-\ell-s}.
\end{equation}
In the radial integral the differences insert
$(|Q|^2-1)^\ell(r^4-1)^s$.  For the angular Fourier coefficient,
shift $\theta$ by a sufficiently small fixed imaginary amount in the
direction giving decay in $|j-k|$.  The continued expression
$Q(re^{i\theta})Q(re^{-i\theta})$ equals one at $r=1$.
Its positive radial logarithmic derivative stays positive under this
small shift, and compactness away from the boundary gives a strict
bound below one there.  It therefore decays as $e^{-c(1-r)}$ and
differs from one by $O(1-r)$.  Radial integration proves
\eqref{c:gram-differences}.

The inverse matrix satisfies
\begin{equation}\label{c:inverse-gram-differences}
 |\Delta^\ell\delta^sR_{jk}|
 \le C_{\ell,s}e^{-c|j-k|}(N+j+k)N^{-\ell}
                   (1+\min\{j,k\})^{-s}.
\end{equation}
First take $\ell=s=0$.  We use the exponential-conjugation argument
for localized inverses; compare Jaffard~\cite[Proposition 2]{Jaffard1990}.
The parameter differences will be retained explicitly below.
Conjugating $P$ by diagonal factors
$\sqrt{N+j}$ gives a positive operator with bounded positive inverse,
by \eqref{c:coefficient-norm}.  Conjugation of this normalized operator
by weights of sufficiently small exponential logarithmic slope changes
it by an arbitrarily small operator norm, by row and column sums in
\eqref{c:gram-differences}.  Its inverse remains bounded.  Truncate
these weights and pass to the limit to obtain exponential decay of
inverse entries, proving the undifferenced case.

Here are details for the differences.  Write
$(SA)_{jk}=A_{j+1,k+1}$, so $\delta A=SA-A$.  The half-line matrix
product has the exact shift rule
\[
 S(AB)=(SA)(SB)+E(A,B),\qquad
 E(A,B)_{jk}=A_{j+1,0}B_{0,k+1}.
\]
In particular $PR=I$ gives, with the $N$ dependence explicit in the
second identity,
\[
 \delta R=-R(\delta P)SR-R E(P,R),
 \qquad
 \Delta R_N=-R_N(\Delta P_N)R_{N+1}.
\]
The first identity displays the boundary contribution that is absent
from a difference in $N$.  Bounds of the form in
\eqref{c:inverse-gram-differences}, with $(N+j+k)^a$ in place of its
first power, compose under matrix multiplication by adding the orders
$a$; convolution of exponential decays absorbs every fixed polynomial
in index distance.  The boundary term $E(P,R)$ and its fixed index
shifts have exponential decay in $j+k$, which pays every displayed
index-decay factor.  First establish all pure $N$ differences from the
second identity.  Then induct on $s$ in the first identity, treating
every fixed $\ell$ at each step.  Discrete product rules use only
lower-order index differences of $R$; the extra difference on $P$ and
the boundary localization give the additional index decay.  Together
with \eqref{c:gram-differences}, these bounds prove
\eqref{c:inverse-gram-differences}.

These estimates also give convergence at the boundary off the even
diagonal.  On each matrix diagonal $j-k$ fixed, put
$n=\min\{j,k\}$ and use a smooth dyadic partition in $n$.
The variable factor is $(v^2\overline u^2)^n$.
Summing by parts $s$ times, and then summing the exponentially decaying
matrix diagonals, bounds the block $n\asymp m$ of
$N^{\ell-1}\Delta^\ell K_N^0$ by $C_sm^{2-s}$, uniformly whenever
$|1-v^2\overline u^2|$ is bounded below.  Choosing $s>2$ gives
uniform convergence of these dyadic blocks and the asserted difference
bound and continuity.  In the open bidisk this sum agrees with the
absolutely convergent kernel, so it gives its boundary values.

For fixed $j,k$, boundary radial Laplace integration shows that
$NP_{jk}$ tends to a positive Toeplitz matrix.  Its symbol is a
positive constant times the reciprocal boundary phase speed of $Q$,
restricted to even powers.  Exponential matrix decay allows passage
to the limit in both inverse identities, so $R_{jk}/N$ tends to the
inverse of this Toeplitz form.  Factor the positive symbol as
$1/|O|^2$, by harmonically extending its logarithm and taking a
harmonic conjugate, normalized by $O(0)>0$.  This is the usual outer
factor construction; see Simon~\cite[Section 5]{Simon2005Szego}.
Real analyticity,
evenness, and reflection symmetry of the symbol give the asserted
analytic continuation, nonvanishing, and symmetries of $O$.
In particular $O$ is real and cannot change sign on either diameter.
Multiplication by $1/O$ is an isometry from the resulting weighted
even Hardy space onto the ordinary even Hardy space.  The latter has
kernel $(1-v^2\overline u^2)^{-1}$.  This identifies the limit;
the same summable dyadic tail bound makes the convergence uniform on
the stated compact sets.
\end{proof}

The model has now isolated the only large part of the evaluation
kernel.  The remaining change of measure contributes a bounded vector,
which is small enough for every residual estimate below.

Let $K_N$ denote the kernel for $\mathcal G_N$.  If
$u=v(t)$, $0\le t\le\lambda$, or $t\uparrow b/2$, then
\begin{equation}\label{c:kernel-correction}
 \|\Delta^\ell(K_N(\cdot,u)-K_N^0(\cdot,u))\|_N
 \le C_\ell N^{-\ell}.
\end{equation}
At $t=b/2$, where $u=1$, the difference of columns has a limit in
$\|\cdot\|_N$.  Neither full boundary column is asserted to be a
finite-norm vector.  We define $K_N(\cdot,1)$ as the model boundary
column plus this finite-norm correction, with evaluations in the open
disk obtained as limits from inside.
Indeed the pairs formed by a point on either fixed interior trace and
these $u$ are uniformly off the even diagonal.  On the traces,
\[
 |J|\asymp e^{-\pi|y|/b},\qquad
 1-|Q|\gtrsim e^{-\pi|y|/b}.
\]
Lemma~\ref{c:model-kernel} and discrete product rules therefore give
\[
 \|\Delta^\ell(T_NK_N^0(\cdot,u))\|_{L^2(\{\pm d_{\mathrm{tr}}+iy\})}
 \le C_\ell N^{-\ell}.
\]
For example the undifferenced integrand is bounded by
\[
 CN e^{-\pi|y|/b}\exp\{-cNe^{-\pi|y|/b}\}.
\]
Its squared integral is bounded; each difference supplies its stated additional
factor.  Apply the bounded correction form to this column and solve
with $\mathcal G_N^{-1}$.  The resolvent identity and the inverse
difference bounds prove \eqref{c:kernel-correction}.  Dominated
convergence on the traces gives its endpoint interpretation.

Evaluation at a fixed interior disk point costs $O(\sqrt N)$ in
\eqref{c:coefficient-norm}.  Consequently, for $v,u$ corresponding to
$0\le t\le\lambda$, the actual kernel has the same strictly positive
leading limit after division by $N$, with $\ell$th differences
$O(N^{1-\ell})$.  The same holds when exactly one variable is $1$.

\subsection{Residual estimates at the two evaluation points}

We prove the following bounds, which are the full analytic input to the
quotient in \eqref{c:strip-residual-formula}:
\begin{equation}\label{c:residual-estimates}
 \begin{split}
 |\Delta^\ell r_Nh_j(-a)|
     &\le N^{-1+b\sqrt a/\pi-\ell+o(1)}\quad(j=0,1),\\
 |\Delta^\ell r_Nc_3(-a)|
     &\le N^{b\sqrt a/\pi-\ell+o(1)}\quad(a=0,4),\\
 -r_Nc_3(0)&\ge c>0.
 \end{split}
\end{equation}
All orders $\ell$ are fixed.  The estimates with $o(1)$ mean that
each fixed positive exponent slack is allowed, with its own constant.

First take the orthogonal residual against $M_N$ alone.  For input
$h=h_j$ its transform is $T_NA_{h,N}$, where $A_{h,N}$ represents the
functional obtained by pairing $h$ with the residual space.  Write
\[
 h(z^2)=\int_{\mathbb R}f_h(y)\cos(yz)\,dy.
\]
The function $f_h$ is integrable against $e^{u|y|}$ for each
$u<\pi/(3\lambda)=\pi/b$.  This follows by Fourier inversion and
shifting within closed narrower strips: the first zeros of
$\mathcal D(z^2)$ have imaginary part $\pi/(3\lambda)$, and the
functions decay exponentially along every such horizontal line.
Thus the functional on $A$ is
$\int f_h(y)(T_NA)(iy)\,dy$.

On $\|A\|_N\le1$, the area mean-value bound and comparison to
$T_{\lfloor N/2\rfloor}$ give
\[
 |(\Delta^\ell T_NA)(iy)|
 \le C_\ell N^{-\ell}
       \exp\{-cNe^{-\pi|y|/b}\}.
\]
Indeed $1-|Q(v(iy))|\gtrsim e^{-\pi|y|/b}$ and
$|1-Q|\lesssim1-|Q|$ there.  Integrating against $f_h$, then using
the inverse-form differences, yields
\begin{equation}\label{c:h-vector-bound}
 \|\Delta^\ell A_{h,N}\|_N\le N^{-1-\ell+o(1)}.
\end{equation}
For an input $C_s$ with $0\le s<b/2$, the representing vector is
instead $J(s)Q(v(s))^NK_N(\cdot,v(s))$.  As $s\uparrow b/2$, all
finite projection pairings and the residual transforms on compact
subsets of the strip converge.  Thus the formula extends to $c_3$
using the boundary column just constructed and $Q(1)=1$; it does not
assert a bounded Hilbert-space evaluation functional at $u=1$.

We now impose the one remaining oblique constraint.  Put
$u_0=v(\lambda/2)$.  The transform of $r_Nh$ is $T_N$ applied to
\begin{equation}\label{c:oblique-vector}
 A_{h,N}-K_N(\cdot,u_0)
       \frac{\chi_N(A_{h,N})}{\chi_N(K_N(\cdot,u_0))},
\end{equation}
where $\chi_N$ averages values at $v(t)$, $0\le t\le\lambda$,
with probability proportional to $J(t)|Q(v(t))|^Ndt$.
To check the formula, orthogonality to $M_N$ is already encoded in
$T_N$, and $\int_0^\lambda C_tdt=g$.  On this interval $Q$ has
constant sign away from its zeros, so this last condition is precisely
$\chi_N=0$.

Every fixed $\ell$th difference of the probabilities defining
$\chi_N$ is $O_\ell(N^{-\ell})$ in total variation.  Divide their
exponential base by its maximum.  Differencing inserts powers of the
gap to one.  Integrals at exponents $N/2$ and $N$ are comparable:
the gap vanishes to finite order at its finitely many maxima, by
analyticity, and the one-dimensional radial Laplace estimate applies
in disjoint neighborhoods of them.  Normalization and the quotient
difference rule give the claim.  Kernel positivity and its leading
limit now show
\[
 \chi_N(K_N(\cdot,u_0))\ge cN,
 \qquad
 |\Delta^\ell\chi_N(K_N(\cdot,u_0))|
 \le C_\ell N^{1-\ell}.
\]
For $h_j$, the subtracted vector in \eqref{c:oblique-vector} has the
same bound as \eqref{c:h-vector-bound}: use the $O(\sqrt N)$ interior
evaluation bound, the $O(\sqrt N)$ column norm, and the denominator
of order $N$, with their corresponding difference bounds.  For $c_3$,
the scalar multiplier of the interior column has $\ell$th differences
$O_\ell(N^{-\ell})$.  The model columns at $1$ and $u_0$ have
$T_N$-images on the imaginary axis bounded by $O(N^{-\ell})$
after $\ell$ differences, by the off-diagonal kernel estimate; their
finite-norm corrections satisfy the same bound.  We have proved,
with $F_{h,N}(\omega)=\int(r_Nh)C_\omega\,d\mu$, that
\begin{equation}\label{c:transform-uniform}
 |\Delta^\ell F_{h,N}(iy)|\le N^{a_h-\ell+o(1)},\qquad
 a_{h_0}=a_{h_1}=-1,\quad a_{c_3}=0,
\end{equation}
uniformly for $y\in\mathbb R$.

\paragraph{Growth on a shifted contour.}
The bound at $-4$ requires a contour argument; it is not an evaluation
on the integration axis.  For $h=h_0,h_1,c_3$, put
\[
 \rho_{h,N}(z)
 =(r_Nh)(z^2)\frac{z\mathcal D(z^2)}{\sinh(4\lambda z)}.
\]
We first prove that, for every fixed $Y>0$ which is not an even integer,
\begin{equation}\label{c:contour-growth}
 \int_{\mathbb R}|\rho_{h,N}(z+iY)|\,dz
 \le C_Y N^{bY/\pi+C},
\end{equation}
where the additive exponent $C$ is independent of $Y$.  This uniformity
will allow us to choose $Y$ after fixing the desired difference order.
The factors $\mathcal D$ remove the poles of $h_0,h_1$ in this product.

In the orthogonal polynomial expansion
\eqref{c:polynomial-generator} with $c=b$,
coefficients of an element of $F_{0,N}$ beyond index $N-1$ have the
form
\[
 q^jp(j)+d_0q_1^j,\qquad \deg p<N,
 \quad 0<q_1=w_{\lambda/2}<q=w_\lambda<1.
\]
Divide by $q^j$ and take $N$ consecutive forward differences to bound
$d_0$ by $e^{O(N)}$ on the unit sphere.  Lagrange interpolation on
$N$ consecutive indices bounds the remaining polynomial.  Consequently
the coefficients decay exponentially beyond $C'N$, relative to the
norm, for a fixed sufficiently large $C'$.
Fix $2\lambda<c<b$ and repeat the coefficient-tail argument in the
orthonormal polynomial basis for $\mu_c$. For each $f\in F_{0,N}$,
its truncation at degree $O(N)$ approximates $f$ with exponentially
small relative $\mu_c$-norm error. The small-circle generating-function
bound gives $|P_j^{(c)}(z^2)|\le C_c^{j+1}e^{cz/4}$ for this basis.
Choosing a fixed $C'$ sufficiently large therefore makes the
truncation's $\mu_c$-norm on $z>C'N$ exponentially small relative to
$\|f\|_{\mu_c}$. By the triangle inequality, a fixed positive fraction
of the full $\mu_c$-norm of $f$ lies on $z\le C'N$. On this interval
the density ratio $d\mu_c/d\mu_b$ is at most $e^{O(N)}$, and hence
\[
 \|f\|_{\mu_c}\le e^{O(N)}\|f\|_{\mu_b},\qquad f\in F_{0,N}.
\]
All constants here may depend on the fixed $c$, not on $N$.

Pairing against $c_3$ with $\mu$ has norm at most a fixed power of
$N$ on this space.  On $z\le N^2$ use Cauchy--Schwarz; since
$c_3(z^2)^2d\mu/dz=O(1+z)$, its truncated squared integral is
polynomial in $N$.  Beyond $N^2$ use the just established $\mu_c$
comparison: the other factor in Cauchy--Schwarz decays as
$e^{-(b-c)z}$, which dominates $e^{O(N)}$.  Thus the orthogonal
projection of $c_3$ on $M_N$ has polynomially bounded norm.  Its
extra oblique projection term also does, by
\eqref{c:oblique-vector}; for $h_j$ this already follows from
\eqref{c:h-vector-bound}.

Finally the generating formula on
$|w_t|=\exp(-1/(j+1))$ gives
\[
 |P_j((z+iY)^2)|
 \le C_Y(j+1)^{1/2+bY/\pi}e^{b|z|/2}.
\]
Here the corresponding parameter satisfies
$|\operatorname{Re}t|<b/2$ and
$|\operatorname{Im}t|\le(b/\pi)\log(j+1)+O(1)$.
Summing with the coefficient bounds, and including the elementary
direct inputs, proves \eqref{c:contour-growth}: multiplication by the
density leaves integrable exponential decay in $|z|$.  In particular
the exponent has coefficient $bY/\pi$ and an additive constant
independent of $Y$, as required.

\paragraph{Recovering the two evaluations.}
Fix the difference order $\ell$ and an arbitrarily small $\eta>0$, and
put $S_N=(b/\pi+\eta)\log N$.  In the full even Fourier integral for
$F_{h,N}(iy)$, shift the $z$ contour to $\operatorname{Im}z=Y$.
The only crossed poles are the simple poles of $1/\sinh(4\lambda z)$
at $z=2in$, $n\ge1$.  The growth bound gives, for $y\ge S_N$,
\begin{equation}\label{c:contour-expansion}
 F_{h,N}(iy)=\sum_{0<2n<Y}b_{n,N}e^{-2ny}
                  +O_Y(e^{-Yy}N^{bY/\pi+C}),
\end{equation}
where $b_{1,N}$ is a fixed nonzero multiple of $r_Nh(-4)$.
Apply \eqref{c:contour-expansion} also at $N+1,\ldots,N+\ell$,
keeping the same $S_N$, and take differences.  At
$y=S_N+j$ for enough consecutive integers $j\ge0$, the exponentials
form a fixed Vandermonde system.  Choose the non-even height $Y>2$
sufficiently large depending on $\ell,\eta$, and then let $N$ tend to
infinity.  Since $C$ is independent of $Y$, the remainder
$N^{-\eta Y+C}$ is negligible at the required difference order.
Equation~\eqref{c:transform-uniform} then bounds every coefficient
$e^{-2nS_N}\Delta^\ell b_{n,N}$ by
$N^{a_h-\ell+o(1)}$.  For $n=1$ this gives
\[
 |\Delta^\ell r_Nh(-4)|
 \le N^{a_h+2b/\pi+2\eta-\ell+o(1)}.
\]
Letting $\eta$ be arbitrarily small proves the claimed exponent.
The same expansion bounds the integral over $y\ge S_N$ by
$N^{a_h-\ell+o(1)}$.  Below $S_N$, integrate
\eqref{c:transform-uniform}; its logarithmic factor is absorbed in
$o(1)$.  Fourier inversion at $z=0$, where the density of $\mu$ is
positive and finite, gives the first two lines of
\eqref{c:residual-estimates} also at $a=0$.

\subsection{A uniform nonzero denominator}

It remains to prove $-r_Nc_3(0)\ge c$.  An upper estimate for the
residual would not suffice to divide by it.
First replace $G_N$ by $G_{N,t}$, $0<t<\lambda$, and use the
comparison measure $W$ in both projection and transform.  Its kernels
are exactly $K_N^0$, and $\chi_N$ becomes point evaluation at $v(t)$.
Apart from the fixed positive factor $J(b/2)$, the residual transform is
$J(iy)Q(v(iy))^N$ times
\[
 K_N^0(v,1)-K_N^0(v,u_0)
          \frac{K_N^0(v(t),1)}{K_N^0(v(t),u_0)},
 \qquad v=v(iy).
\]
After division by $N$, Lemma~\ref{c:model-kernel} gives a limit
uniform on the imaginary diameter.  At $v=\pm i$, its sign is the
sign of
\[
 \frac12-
 \frac{1-v(t)^2u_0^2}{(1+u_0^2)(1-v(t)^2)},
\]
which is strictly negative, uniformly for $0\le t\le\lambda$.
Indeed the second term is at least $1/(1+u_0^2)>1/2$.
The remaining factors $O$ are positive on both diameters, with uniform
positive bounds on the compact sets in question.

Fourier inversion at zero integrates the bracket divided by $N$ against
\[
 cNJ(iy)Q(v(iy))^Ndy,\qquad c>0.
\]
This positive measure has mass bounded above and below and concentrates
at $y\to\pm\infty$: on the imaginary diameter $Q\ge0$,
$1-Q\asymp e^{-\pi|y|/b}$, and
$J\asymp e^{-\pi|y|/b}$.
The boundedness supplied by Lemma~\ref{c:model-kernel} and the negative
endpoint limit therefore give the desired lower bound for $W$,
uniformly in $t$.

To pass to $\mu$, its exact density ratio to $W$ is
\[
 \frac{d\mu}{dW}
 =\frac{(2u-1)^2(2u+1)}{4u(2u^2-1)},
 \qquad u=\cosh(\lambda\sqrt x).
\]
It is positive and smooth on $[0,\infty)$ and eventually increasing.
For $r(x)=(x+2)/(x+1)$, choose a fixed integer $j>0$ so that
$\mu/(r^jW)$ is increasing: on a compact interval the negative part
of the logarithmic derivative is bounded, while
$-(\log r)'$ is strictly positive; outside it the original derivative
is already nonnegative.  The ordered association input then bounds
the expectation of $-e_0c_3$ for $\mu$ below by that for $r^jW$.

Each multiplication of the base weight by $r$ changes this expectation
by only a bounded factor.  Indeed it multiplies it by the ratio of the
characteristic-product expectations of $\prod_i r(x_i)$ under
$(A_N,G_{N,t})$ and $(F_{0,N},G_{N,t})$ at the current base weight.
The interlacing input compares the first to $(U_N,U_N)$ and the second
to $(M_N,M_N)$; symmetric codimension-one compression compares these
last two.  The bounded range of $\log r$ gives a uniform constant at
each step, independent of $N,t$.  After the fixed number $j$ of steps,
the lower bound therefore remains positive.  Finally integrate the
positive $G_{N,t}$ mixture to obtain it for $G_N$.
This proves the last line of \eqref{c:residual-estimates}.

\paragraph{The quotient.}
We can now finish Proposition~\ref{c:strip-smoothness}.
The denominator $r_N(c_3-h_1)(0)$ in
\eqref{c:strip-residual-formula} is bounded away from zero for all
sufficiently large $N$, because
$-r_Nc_3(0)\ge c$ and $|r_Nh_1(0)|\le N^{-1+o(1)}$.
Its reciprocal has $\ell$th differences $N^{-\ell+o(1)}$, by the
discrete inverse rule and \eqref{c:residual-estimates}.  In the
numerator, $r_Nh_0(-4)$ has differences
$N^{-1+2b/\pi-\ell+o(1)}$, while the product of
$r_N(c_3-h_1)(-4)$ and $r_Nh_0(0)$ has exactly the same bound.
Apply the discrete product and quotient rules.  Since
$2b/\pi=3/4$, the resulting exponent is $-1/4-\ell+o(1)$,
as asserted.  For each fixed $\ell$ and exponent slack, the finitely
many remaining values of $N$ are absorbed in $C_{\ell,\eta}$.
